%======================================================================
\subsection{Flexible sub-assemblies with routed items --- Built, 30 Oct}\label{sec:flexsub-built}
%======================================================================

Answer W11 of the flexible-wire plan (\S\ref{sec:rou-flex-plan}) is built: a sub-assembly inserted \emph{flexible} now
brings its route paths, wires, hoses, pipes and bundles into the parent, and they follow the parts the parent moves.

\paragraph{What you can do now.} Insert a sub-assembly with \code{Rigid=False}. Its routed items appear inside the
\code{AssemblyLink} next to the links to its parts. Their attachments are re-pointed at those links, so a wire
between two pins stretches or reshapes when the parent's joints (or you) move one housing. Edit the wire in the
sub-assembly's file (gauge, colour, a new wire, a deleted one) and the parent follows on the next update. Make the
link rigid again and the copies go; make it flexible again and they return.

\paragraph{Decided while building.}
\begin{itemize}[nosep]
  \item \textbf{Copies, not links.} An \code{App::Link} to a wire shows the source's shape at the source's place and
        cannot carry other attachments, so it can never follow a part the parent moves. A copy is an ordinary routed
        object that recomputes in the parent from its own attachments. Each copy has a hidden \code{FlexSource}
        property (the source's name).
  \item \textbf{Split of work.} The C++ side (\code{AssemblyLink.cpp}, \code{[BW]}) stops linking routed items as if
        they were parts (they are \code{Part::FeaturePython} with \code{BendOK}, so the old code made a static
        \code{App::Link} to each), leaves copies alone when it prunes the group, and calls
        \code{bruceworx\_flexsub.sync(link, \{name: link\}, assembly)} after every update. The copying is Python
        (\file{bruceworx\_flexsub.py}), so it follows whatever properties the routing classes grow; without the
        Heuretes FreeCAD module the call is silent.
  \item \textbf{Re-pointing.} An attachment \code{(object, sub)} below a component becomes
        \code{(link to the component, path from the component + sub)}: \code{(HA, "Pin1.")} becomes
        \code{(LHA, "Pin1.")}, a frame inside the part becomes \code{(LHA, "Pin1.")} too. One outside the
        sub-assembly is dropped and the point stays free.
  \item \textbf{Properties.} Editable values are copied when they differ; outputs (read-only, \code{RoutePoints},
        the shape) are the copy's own. A path's free points follow the source, the attached ones are re-read from the
        links. Materials and links are never copied (they would touch the copy at every update).
  \item \textbf{Known edge.} Until the parent's solver or you place the links, both housings sit at the origin and a
        wire between them cannot be swept: one console error at insert, gone at the first placement.
\end{itemize}

\paragraph{Tests.} \file{TestBruceWorxFlexSub.py}, 4 headless tests with the sub-assembly in its own file: copies
made and attachments re-pointed (and the source untouched); moving a housing in the parent moves the copy's end
exactly and reshapes the wire; gauge changes, new and deleted wires are synced; rigid removes and flexible restores;
save, reload, no doubling and still following. Not done: the battery check on a vault copy (the harness there needs
the rope solver, ROU-08, to be worth running).

\paragraph{Open.} Copies are rebuilt from the source; a value typed on a copy is overwritten at the next update. Say
if you want a per-copy override list.

\textbf{Answer (30 Oct):} copies are read-only; edit the wire in its sub-assembly.
